\section{Training Behavior of the Compared Approaches}

\subsection{Expected behavioral signatures}
\begin{draftpoints}
  \item Fixed CL without replay should acquire the incoming class through global parameter updates but is expected to overwrite earlier representations because no historical signal constrains those updates
  \item NDL without replay protects mature encoder features during local plasticity and adds capacity, so it may reduce forgetting even if its local objective yields worse total reconstruction
  \item Original-data replay provides the strongest stabilization signal and therefore separates growth quality from replay-generation quality
  \item Intrinsic replay can preserve old representations only if its latent Gaussian samples decode into sufficiently faithful old-class images; poor replay can simultaneously distort stabilization and the thresholds used to judge novelty
  \item method that reports low forgetting but high incoming-class error is stable but insufficiently plastic; both acquisition and retention must therefore be shown
\end{draftpoints}

\subsection{Observed MNIST behavior}
\begin{draftpoints}
  \item CL with original-data replay has the best confirmatory reconstruction endpoint and nearly zero positive forgetting, showing that a fixed network of the final capacity can learn the curriculum when clean old-class samples remain available
  \item NDL with original-data replay also retains old classes well, but its macro MSE is more than three times the matched CL value; local growth and stabilization do not convert the added parameters into equally useful global reconstructions
  \item NDL without replay forgets less than matched CL without replay (0.00187 versus 0.00897 positive forgetting), but its macro MSE is 1.91 times larger; growth yields a real retention trade-off rather than a complete continual-learning improvement
  \item NDL with intrinsic replay updates every incoming class, reaches the same large endpoint $[232,140,91,44]$ in every seed, and exhausts every class/depth growth loop without reaching the accepted outlier quota
  \item intrinsic condition requires about 88,292 incremental model updates on average, compared with 1,113 for its matched CL control, yet ends with higher reconstruction error and higher forgetting
\end{draftpoints}

\begin{figure}[htbp]
\centering
\includegraphics[width=\textwidth]{mechanism_comparison.png}
\caption{Mechanism-level comparison of endpoint capacity, growth behavior, and performance. The figure combines architecture and diagnostic outcomes so that a paper-like final width is not interpreted independently of cap exhaustion and reconstruction quality. Source tables: \artifact{mnist_architecture.csv}, \artifact{mnist_condition_metrics.csv}, and \artifact{mnist_seed_metrics.csv}.}
\label{fig:mechanism-comparison}
\end{figure}

\subsection{Growth-stop behavior}
\begin{draftpoints}
  \item Quota-stop fraction is the share of growth loops that stop because error demand has been satisfied; a value near one supports organic growth, while a value near zero indicates that the architecture is limited by an allowance
  \item Clean-replay and no-replay NDL reach a mean quota-stop fraction of 0.875 because growth occurs in a minority of class-level combinations and most evaluated loops require no update
  \item Intrinsic-replay NDL has quota-stop fraction 0.000, updated-class fraction 1.000, and 32 unresolved exhausted level loops in every seed, indicating systematic rather than isolated novelty pressure
  \item fixed endpoint under intrinsic replay is therefore not evidence that the same useful architecture emerged repeatedly; it is evidence that the configured allowances bound the same repeated failure mode
\end{draftpoints}

\subsection{SD-19 behavior}
\begin{draftpoints}
  \item Both SD-19 conditions end at $[1100,600,350,70]$, so the comparison holds parameter count constant at 3,518,754
  \item NDL uses roughly 7,147--7,432 model updates compared with 72 for CL and still has about twice the macro and foreground reconstruction error
  \item NDL stops only 12.5\% of loops by quota and retains 21 unresolved exhausted levels per seed, showing that greater task complexity does not resolve the cap-driven behavior
\end{draftpoints}

\subsection{Trajectory data required for a complete behavioral account}
\begin{draftpoints}
  \item preserved export supports seed endpoints, per-class final errors, and detailed growth reports for selected conditions, but it does not contain a uniform cross-condition time series after every curriculum step
  \item future matched rerun should export one row per condition, seed, class step, and depth with incoming MSE, old-class MSE, widths before and after growth, nodes added, threshold, initial and final outliers, stop reason, updates, and runtime
  \item \missingartifact{A complete \artifact{training_trajectory.csv} is needed before drawing class-step curves for all six confirmatory conditions; no trajectory curve is inferred from endpoint tables.}
\end{draftpoints}

\begin{table}[htbp]
\centering
\caption{Schema proposed for the missing uniform training-trajectory export.}
\label{tab:trajectory-schema}
\scriptsize
\begin{tabularx}{\textwidth}{p{2.8cm}Y}
\toprule
Field group & Required columns \\
\midrule
Identity & condition, seed, class step, class ID, autoencoder depth \\
Performance & macro MSE, incoming-class MSE, old-class MSE, positive forgetting \\
Growth & width before, width after, nodes requested, nodes accepted, growth rounds \\
Outliers & threshold, initial outliers, final outliers, allowed outliers, stop reason \\
Cost & plasticity updates, stability updates, next-layer updates, cumulative runtime \\
\bottomrule
\end{tabularx}
\end{table}
